\usepackage{ctex}
\usepackage{multicol}
\usepackage{geometry}
\geometry{left=0.5in, right=0.5in, top=0.5in, bottom=0.5in}
\usepackage{titlesec}
\usepackage{paralist}
\usepackage{enumitem}
\usepackage{amssymb}

\usepackage{setspace}
\usepackage{geometry}%调用页面元素尺寸设置宏包
\usepackage{fancyhdr}%调用页面样式设置宏包
\pagestyle{fancy}%使用fancy页面样式，与上一个命令紧跟着
\usepackage{lastpage}%调用末页标签宏包
\usepackage{xcolor}
\usepackage{soul}
\usepackage{colortbl}
\usepackage{dashbox}
\usepackage[utf8]{inputenc}
\usepackage[T1]{fontenc}
\usepackage{ruby}  % 引入 ruby 宏包
\usepackage{tikz}  % 用于绘制图形
%\usepackage{tcolorbox}
\usepackage{titlesec}
\usepackage{booktabs}  % 用于创建漂亮的表格，支持三线表等
\usepackage{setspace}  % 调整文档行距
\usepackage{amsmath}   % 提供数学公式排版支持
\usepackage{stfloats}  % 允许在双栏文档中使用跨页的浮动对象
\usepackage{adjustbox}
\usepackage{graphicx}  % 支持插入图片
\usepackage{datetime}  % 提供日期和时间格式化
\usepackage{fancyhdr}  % 自定义页眉页脚 
\usepackage{caption}   % 改进浮动体标题的格式,提供\captionof命令
\usepackage{array}     % 提供更多的表格列格式选项
\usepackage{float}     % 提供浮动体排版控制
\usepackage{makecell}  % 允许表格中的单元格内换行
\usepackage{hyperref}  % 创建超链接，使文档可交互
\usepackage{multicol}  % 允许文档分为多栏
\usepackage{titlesec}  % 用于自定义章节标题的样式
\usepackage[ruled,linesnumbered]{algorithm2e}  % 引入算法环境宏包
\usepackage{ifthen}    %这个宏包提供逻辑判断命令
\usepackage{bm} %粗体数学符号的支持
\usepackage{marginnote}
\usepackage{footnote}
\usepackage{mdframed}
\usepackage{multirow}
\setlength{\parindent}{0pt}
\usepackage{tikz}  % 用于绘制图形
\usetikzlibrary{shapes.geometric,calc,arrows,shadows.blur}

% 定义新的列类型用于自动换行
\newcolumntype{P}[1]{>{\raggedright\arraybackslash}p{#1}}
\usetikzlibrary{positioning}
%设置主要字体为TimesNew Roman
\setmainfont{Times New Roman} 
% 为章节标题定义新字体 Times New Roman
\newfontfamily\sectionef{Times New Roman} 
% 为中文章节标题定义新的楷书字体
\newCJKfontfamily\sectioncf{heiti}
\usepackage{makeidx} % 用于索引
\usepackage{tocloft}        % 用于自定义目录格式
\usepackage{hyperref}       % 用于创建可点击的链接
\usepackage{xeCJK}         % 用于中文支持
\usepackage{tcolorbox} 
\tcbuselibrary{most}
% 设置tcolorbox的样式
\usepackage{tikz}
\usepackage{forest}



%颜色定义，c0-c4：由深到浅----------------------------
\definecolor{c0}{HTML}{647382}
\definecolor{c1}{HTML}{9AA6B2}
\definecolor{c2}{HTML}{BCCCDC}
\definecolor{c3}{HTML}{D9EAFD}
\definecolor{c4}{HTML}{F8FAFC}

%导入盒子定义-----------------------------------------
%定义知识点拓展盒子knowledgebox----------------------------------
\newtcolorbox{knowledgebox}[2][]
  {enhanced,
    before skip=2mm,
    after skip=2mm,
    breakable,
    colback=c4,
    colframe=c1,%边框颜色
    boxrule=0.5mm,
    attach boxed title to top left={xshift=1cm,yshift*=1mm-\tcboxedtitleheight},
    fonttitle=\bfseries,
    % varwidth boxed title*=-3cm,
	boxed title style={frame code={
                    \path[fill=c1]
                    ([yshift=-1mm,xshift=-1mm]frame.north west)
                    arc[start angle=0,end angle=180,radius=1mm]
                    ([yshift=-1mm,xshift=1mm]frame.north east)
                    arc[start angle=180,end angle=0,radius=1mm];
                    \path[left color=c1,right color=c1,
                        middle color=c1]
                    ([xshift=-2mm]frame.north west) -- ([xshift=2mm]frame.north east)
                    [rounded corners=1mm]-- ([xshift=1mm,yshift=-1mm]frame.north east)
                    -- (frame.south east) -- (frame.south west)
                    -- ([xshift=-1mm,yshift=-1mm]frame.north west)
                    [sharp corners]-- cycle;
                },interior engine=empty,
        },
    title=\large \( \textbf{解释说明： #2} \),#1}

%定义解题点拨盒子skillbox----------------------------------
\newtcolorbox{skillbox}[2][]
  {enhanced,
  breakable,
  sharp corners,
  attach boxed title to top left={
    yshifttext=-1mm
  },
  colback=c4,
  colframe=c1,
  boxrule=0.5mm,
  fonttitle=\bfseries,
  coltitle=white,
  boxed title style={
    sharp corners,
    size=small,
    colback=c1,
    colframe=c1,
  },
  title=\large \( \textbf{解题点拨： #2} \),
  #1}



%定义考查方式盒子pointbox----------------------------------
\newtcolorbox{pointbox}[2][]
  {enhanced,
  breakable,
  colback = c4,
  frame hidden,
  borderline west = {2pt}{0pt}{c1},
  coltitle = c0,
  fonttitle = \bfseries\sffamily,
  boxrule = 0sp,
  sharp corners,
  detach title,
  before upper = \tcbtitle\par\smallskip,
  description font = \mdseries,
  separator sign none,
  segmentation style={solid, c1},
	title=\large \( \textbf{考查方式： #2} \),#1}

% 备用盒子1boxa----------------------------------
\makeatletter
\newtcolorbox{boxa}[1]{%
    enhanced,
    breakable,
    colback=c4,
    colframe=c0,
    attach boxed title to top left={yshift*=-\tcboxedtitleheight},
    fonttitle=\bfseries,
    title={\large \( \textbf{备用盒子： #1} \)},
    boxed title size=title,
    boxed title style={%
            sharp corners,
            rounded corners=northwest,
            colback=c0,
            boxrule=0pt,
        },
    underlay boxed title={%
            \path[fill=c0] (title.south west)--(title.south east)
            to[out=0, in=180] ([xshift=5mm]title.east)--
            (title.center-|frame.east)
            [rounded corners=\kvtcb@arc] |-
            (frame.north) -| cycle;
        }
}
\makeatletter

%例题盒子定义example----------------------------------
\newtcolorbox{example}[1]{
    colback = c4,
    breakable,
    colframe = c1,
    coltitle = c0,
    title={\large \( \textbf{Example： #1} \)},
    boxrule = 1pt,
    sharp corners,
    detach title,
    before upper=\tcbtitle\par\smallskip,
    fonttitle = \bfseries,
    description font = \mdseries
}

%定义注释盒子note----------------------------------
\newtcolorbox{note}[1][]{%
    enhanced jigsaw,
    colback=c4,
    colframe=c1,
    size=small,
    boxrule=1pt,
    title=\textbf{Note},
    halign title=flush center,
    coltitle=black,
    breakable,
    drop shadow=c2,
    attach boxed title to top left={xshift=1cm,yshift=-\tcboxedtitleheight/2,yshifttext=-\tcboxedtitleheight/2},
    minipage boxed title=1.5cm,
    boxed title style={%
        colback=white,
        size=fbox,
        boxrule=1pt,
        boxsep=2pt,
        underlay={%
            \coordinate (dotA) at ($(interior.west) + (-0.5pt,0)$);
            \coordinate (dotB) at ($(interior.east) + (0.5pt,0)$);
            \begin{scope}
                \clip (interior.north west) rectangle ([xshift=3ex]interior.east);
                \filldraw [white, blur shadow={shadow opacity=60, shadow yshift=-.75ex}, rounded corners=2pt] (interior.north west) rectangle (interior.south east);
            \end{scope}
            \begin{scope}[c1]
                \fill (dotA) circle (2pt);
                \fill (dotB) circle (2pt);
            \end{scope}
        },
        },
        #1,
}

%定义考点习题标题及分割线\et----------------------------------
\newcommand{\et}[0]{
  \color{c1}
  \subsubsection*{\large 考点习题}
  \hrule
  \vspace{1em}
  \color{black}
}

%初始盒子1----------------------------------------
\newtcolorbox{warningbox}{
	colback=green!10,
	colframe=olive!50!black,
	leftrule=0.5mm,
	rightrule=0mm,
	toprule=0mm,
	bottomrule=0mm,
	arc=2pt,
	outer arc=2pt
}

% 初始盒子2：定义定理框样式------------------------------
\newtcolorbox{theorembox}{
	colback=gray!10,
	colframe=olive!50!black,
    leftrule=0.5mm,
    rightrule=0mm,
    toprule=0mm,
    bottomrule=0mm,
	arc=2pt,
	outer arc=2pt
}
%----------------------------------------------------
\newcommand{\view}[3][评论]{%
	\begin{compactdesc}
		\item [#1\ifx\relax#2\relax\else (#2)\fi.] #3
	\end{compactdesc}
}


\usetikzlibrary{mindmap,trees,arrows.meta}
\usetikzlibrary{shapes.misc}  % 为了使用 cross out 样式
\tcbset{
	%enhanced,
	boxrule=0pt, % 边框宽度
	left=3mm, right=3mm, top=3mm, bottom=3mm, % 内边距
	%	width=\textwidth
}

%启用下面三行，可将图、表、参考文献的中文题注和标题变为英文
\renewcommand{\refname}{References}
\captionsetup[figure]{name=图,labelsep=colon}
\captionsetup[table]{name=表,labelsep=colon}

\setlist[itemize]{itemsep=\parsep,parsep=\parskip,topsep=\parskip}
\setlist[enumerate]{itemsep=\parsep,parsep=\parskip,topsep=\parskip}

%习题----------------------------------------------------------
% 定义习题环境
\newcounter{exercise}
\newenvironment{exercise}[1][]{\refstepcounter{exercise}\par\medskip
	\noindent 
  \textbf{习题\theexercise. #1} \rmfamily\nopagebreak[4]}{\medskip}

% 定义选项环境
\newenvironment{choices}{
	\begin{enumerate}[label=\Alph*.]
	}{\end{enumerate}}

% 定义答案环境
\newenvironment{answer}{
	\noindent\textbf{答案：}}{}

% 定义解析环境
\newenvironment{solution}{
	\noindent\textbf{解析：}}{}

%习题----------------------------------------------------------


\hypersetup{
	colorlinks=false,     % 不使用颜色
	pdfborder={0 0 0},   % 设置边框宽度为0
	hidelinks            % 完全隐藏链接标记
}


\renewcommand{\contentsname}{\huge 考点目录}

\makeatletter
% 定义新命令用于标记PE1题目
\newcommand{\peOneNumber}[1]{%
	{\normalsize \colorbox{black!70}{\textcolor{white}{\textbf{PE1.#1}}}}%    % 在题目旁显示编号
	\addtocontents{peOne}{PE1.#1 \dotfill \thepage\par}%
}

% 生成索引的命令
\newcommand{\listofpeOne}{%
	\thispagestyle{plain}%    % 移除当前页的页眉页脚
	{\huge \textbf{PE1题目索引}}\par
	\@starttoc{peOne}%
}
\makeatother

\makeatletter
% 定义新命令用于标记PE1题目
\newcommand{\peTwoNumber}[1]{%
	{\normalsize \colorbox{black!70}{\textcolor{white}{\textbf{PE2.#1}}}}%    % 在题目旁显示编号
	\addtocontents{peTwo}{PE2.#1 \dotfill \thepage\par}%
}

% 生成索引的命令
\newcommand{\listofpeTwo}{%
	\thispagestyle{plain}%    % 移除当前页的页眉页脚
	{\huge \textbf{PE2题目索引}}\par
	\@starttoc{peTwo}%
}
\makeatother

\makeatletter
% 定义新命令用于标记PE1题目
\newcommand{\peThreeNumber}[1]{%
	{\normalsize \colorbox{black!70}{\textcolor{white}{\textbf{PE3.#1}}}}%    % 在题目旁显示编号
	\addtocontents{peThree}{PE3.#1 \dotfill \thepage\par}%
}

% 生成索引的命令
\newcommand{\listofpeThree}{%
	\thispagestyle{plain}%    % 移除当前页的页眉页脚
	{\huge \textbf{PE3题目索引}}\par
	\@starttoc{peThree}%
}
\makeatother

\usepackage{fix-cm}  % 允许任意大小的Computer Modern字体
\renewcommand{\abstractname}{说明}
% 在导言区添加以下命令来重新定义 subsection 的格式
\renewcommand{\thesubsection}{考点\arabic{subsection}}
\renewcommand{\thesubsubsection}{\arabic{subsubsection}}
\usepackage{tocloft}
\setlength{\cftsecnumwidth}{2em}  % section编号的宽度
\setlength{\cftsubsecnumwidth}{4em}  % subsection编号的宽度
\setlength{\cftsubsubsecnumwidth}{2em}  % subsubsection编号的宽度


